\section{The MLIPs Ontology}
\label{sec:ontology}
%% ====================================================================

The \MLIPsOntology{} is an OWL\,2~DL ontology organized into three modules
that address the competency questions identified in
Section~\ref{sec:requirements}. We first introduce the ontology through a
concrete example (Section~\ref{sec:running-example}), then describe each
module formally (Sections~\ref{sec:algorithm-module}--\ref{sec:benchmark-module}).
Figure~\ref{fig:ontology-overview} provides an overview. The ontology
includes 27 formal OWL axioms (12 existential restrictions, 6 cardinality
constraints, 5 inverse existence axioms, 2 property chains, and 2 domain
axioms) that enforce data completeness and consistency; the full axiom
catalog with examples is published alongside the ontology at
\url{\mlipsns}. The ontology uses the namespace \texttt{\mlipsns} (prefix
\mlipsprefix{}); Table~\ref{tab:prefixes} lists the external prefixes used.
In the following, we omit the \mlipsprefix{} prefix for concepts and
properties of the \MLIPsOntology{}.

\subsection{Running Example: Moment Tensor Potentials for Ti-Al}
\label{sec:running-example}

We introduce the ontology through a concrete example: training a Moment
Tensor Potential (MTP)~\cite{shapeev2016mtp,novikov2021mlip} for the
titanium-aluminium (Ti-Al) binary alloy system. This example illustrates
how the ontology captures the full lifecycle of an MLIP study---from the
physics of the simulation cell, through the algorithm's hyperparameters
and training data, to the evaluation of the trained model.

\paragraph{The simulation problem.}
Consider a simulation cell containing $n$ atoms at positions
$\mathbf{r}_1, \ldots, \mathbf{r}_n$ with atomic types
$z_1, \ldots, z_n$ (where $z_i \in \{\text{Ti}, \text{Al}\}$ for our
Ti-Al system). The goal of an interatomic potential is to predict the
total energy $E$ and the forces $\mathbf{F}_i = -\nabla_{\mathbf{r}_i} E$
acting on each atom, given only the atomic positions and types.
Figure~\ref{fig:simulation-cell} illustrates the simulation cell and the
concept of atomic neighborhoods.

\begin{figure}[t]
\centering
\begin{subfigure}[t]{0.48\textwidth}
  \centering
  %% TODO: create 3D simulation cell figure
  \fbox{\parbox{0.9\linewidth}{\centering\vspace{3cm}%
    3D simulation cell with Ti (blue)\\and Al (orange) atoms%
    \vspace{3cm}}}
  \caption{A Ti-Al simulation cell with periodic boundary conditions.}
  \label{fig:simulation-cell-3d}
\end{subfigure}
\hfill
\begin{subfigure}[t]{0.48\textwidth}
  \centering
  %% TODO: create 2D cutoff radius diagram
  \fbox{\parbox{0.9\linewidth}{\centering\vspace{3cm}%
    2D cross-section showing $R_{\text{cut}}$\\and neighborhood $\neigh{i}$%
    \vspace{3cm}}}
  \caption{The neighborhood $\neigh{i}$ of atom~$i$: all atoms within
    $R_{\text{cut}}$ contribute to $V(\neigh{i})$. Atoms beyond
    $R_{\text{cut}}$ (grayed out) are excluded.}
  \label{fig:simulation-cell-cutoff}
\end{subfigure}
\caption{(a)~A Ti-Al simulation cell. (b)~The cutoff radius
  $R_{\text{cut}}$ defines the neighborhood $\neigh{i}$, which
  includes relative positions $\mathbf{r}_{ij}$ and atomic types $z_j$
  of the neighbors.}
\label{fig:simulation-cell}
\end{figure}

In quantum mechanics, the total energy can be computed exactly (within
the approximations of density functional theory) by solving the
Kohn--Sham equations~\cite{kohn1965dft}. However, this is
computationally expensive---scaling as $O(n^3)$ with the number of
atoms---limiting DFT to systems of a few hundred atoms. Machine learning
interatomic potentials approximate the DFT energy surface at a fraction
of the cost, enabling simulations of millions of atoms.

\paragraph{The MTP algorithm.}
In the MTP framework~\cite{shapeev2016mtp}, the total energy of an
atomic configuration $\text{cfg}$ is decomposed into local contributions:
%
\begin{equation}\label{eq:mtp-energy}
  E^{\text{mtp}}(\text{cfg}) = \sum_{i=1}^{n} V(\neigh{i}),
\end{equation}
%
where $\neigh{i}$ is the \emph{neighborhood} of atom~$i$, defined as
the set of all atoms within the cutoff radius $R_{\text{cut}}$. The
neighborhood consists of the atomic type $z_i$ of the central atom, and
for each neighbor $j$ within $R_{\text{cut}}$: the atomic type $z_j$ and
the relative position vector $\mathbf{r}_{ij} = \mathbf{r}_j -
\mathbf{r}_i$. Crucially, atoms beyond $R_{\text{cut}}$ do not
contribute to $V(\neigh{i})$---this \emph{locality assumption} is
what makes the potential computationally efficient.

The local energy $V(\neigh{i})$ is expressed as a linear combination
of basis functions:
%
\begin{equation}\label{eq:mtp-basis}
  V(\neigh{i}) = \sum_{\alpha} \xi_\alpha \, B_\alpha(\neigh{i}),
\end{equation}
%
where $B_\alpha$ are the \emph{moment tensor descriptors}---invariant
functions of the neighborhood that capture the local atomic
environment---and $\boldsymbol{\xi} = \{\xi_\alpha\}$ are the
\emph{parameters to be learned} by fitting to DFT reference data. The
basis functions $B_\alpha$ are constructed from contractions of moment
tensors $M_{\mu,\nu}$ involving radial functions $Q^{(\beta)}(|\mathbf{r}_{ij}|)$
that vanish smoothly at $R_{\text{cut}}$.

\paragraph{Hyperparameters.}
The MTP algorithm has several hyperparameters that control the model's
expressiveness and computational cost:
%
\begin{itemize}
\item $R_{\text{cut}}$ (\emph{cutoff radius}): determines the size of
  atomic neighborhoods. A larger cutoff captures longer-range
  interactions but increases computational cost.
\item $R_{\text{min}}$ (\emph{minimum distance}): the shortest expected
  interatomic distance, used to define the radial basis domain.
\item $N_Q$ (\emph{number of radial basis functions}): controls the
  resolution of the radial description.
\item $\text{lev}_{\text{max}}$ (\emph{maximum level}): determines which
  basis functions $B_\alpha$ are included, controlling the trade-off
  between accuracy and model size.
\end{itemize}
%
These are \emph{not} learned from data---they are chosen by the
researcher before training. In contrast, the coefficients
$\boldsymbol{\xi}$ are determined by fitting
Equation~\eqref{eq:mtp-basis} to DFT-computed energies, forces, and
stresses via linear regression.

\paragraph{Training data.}
The ground truth for training comes from DFT calculations: for a set of
atomic configurations of the Ti-Al system, one computes the total energy,
forces on each atom, and (optionally) stresses. Each configuration is a
snapshot of atomic positions in a periodic cell, generated by sampling
from molecular dynamics trajectories, random perturbations, or active
learning~\cite{novikov2021mlip}. A typical training dataset for Ti-Al
might contain 15{,}000 configurations computed with the PBE
exchange-correlation functional, PAW pseudopotentials, and a plane-wave
energy cutoff of 520~eV using the VASP code.

\paragraph{Ontology mapping.}
This example maps directly onto the \MLIPsOntology{} as follows.
The MTP algorithm is an instance of $\MLIPAlgorithmC$ with four
$\HyperparameterC$ instances ($R_{\text{cut}}$, $R_{\text{min}}$, $N_Q$,
$\text{lev}_{\text{max}}$). The Ti-Al material is a $\MaterialSystemC$.
The 15{,}000 configurations form a $\TrainingDatasetC$, linked to a
$\DFTCalculationC$ with $\DFTSettingsC$ recording the PBE functional and
PAW pseudopotentials.

A $\TrainingRunC$ records the concrete training event: it $\executesR$
the MTP algorithm, $\runsOnR$ the Ti-Al dataset, and $\producesR$ a
$\TrainedModelC$---the specific set of coefficients $\boldsymbol{\xi}$
fitted to this data with $R_{\text{cut}} = 5.0$~\AA{},
$\text{lev}_{\text{max}} = 16$, and $N_Q = 8$. These concrete
values are recorded as $\HyperparameterSettingC$ instances.

When the trained model is evaluated on a test set (e.g., computing the
RMSE of predicted energies against DFT reference values), the
evaluation is captured as a $\BenchmarkResultC$ that $\evaluatesModelR$
the trained model, $\targetMaterialR$ the Ti-Al system, and reports
$\AccuracyMetricC$ instances (e.g., RMSE of energy = 1.2~meV/atom).

Figure~\ref{fig:mtp-ontology-mapping} shows how these ontology instances
relate to each other through the ontology's roles.

\begin{figure}[t]
\centering
%% TODO: create instance diagram showing the MTP/Ti-Al example
\fbox{\parbox{0.9\textwidth}{\centering\vspace{3cm}%
  Instance diagram: MTP algorithm → training run → trained model → benchmark result,\\
  with hyperparameter settings, DFT calculation, and accuracy metrics.%
  \vspace{3cm}}}
\caption{Ontology instance diagram for the running example: training an
  MTP model for Ti-Al. Yellow nodes are instances; blue nodes are classes
  from the \MLIPsOntology{}. The diagram shows how the MTP algorithm,
  training data, trained model, and benchmark results are connected
  through the ontology's roles.}
\label{fig:mtp-ontology-mapping}
\end{figure}


\begin{figure}[t]
\centering
\resizebox{\textwidth}{!}{% Manually-laid-out ontology overview.
% Module membership is encoded by node fill colour:
%   * Method module        → blue   (algofill)
%   * Training Data module → green  (trainfill)
%   * Benchmark module     → orange (benchfill)
%
% Layout invariants:
%   * Six rows defined by the central column (Implementation, MLIPMethod,
%     MLIPRun, TrainedModel, MaterialSystem, TrainingDataset).  All
%     central-column nodes share x = 5, the same minimum width, and a
%     uniform vertical spacing of 1.3 cm.
%   * Left column (Library, SimulationType, AtomicConfiguration) shares
%     west x = -1.1 (anchor=west).
%   * Right column (BenchmarkStudy, BenchmarkResult, AccuracyMetric,
%     MetricProperty) shares east x = 16 (anchor=east).
%   * Four leaf classes (DatasetProvenance, CoveredProperty, MetricType,
%     MetricProperty) sit on a new row at y = -0.8.
%   * Cross-module edges are dashed; the trainedWith and trainedUsing
%     shortcuts are dotted.
%   * Node and edge labels use the same DL macros as the main text
%     (\dlConcept and \dlRole), rendered in math sans-serif.
\begin{tikzpicture}[
    x=1cm, y=1cm,
    every node/.style={font=\footnotesize},
    methodnode/.style={draw, rounded corners=2pt, fill=algofill,
        minimum height=0.55cm, inner sep=3pt, align=center,
        font=\footnotesize, text height=1.6ex, text depth=0.4ex},
    trainnode/.style={draw, rounded corners=2pt, fill=trainfill,
        minimum height=0.55cm, inner sep=3pt, align=center,
        font=\footnotesize, text height=1.6ex, text depth=0.4ex},
    benchnode/.style={draw, rounded corners=2pt, fill=benchfill,
        minimum height=0.55cm, inner sep=3pt, align=center,
        font=\footnotesize, text height=1.6ex, text depth=0.4ex},
    centerwidth/.style={minimum width=2.5cm},
    edge/.style={->, >=stealth, thin},
    crossedge/.style={->, >=stealth, thin, dashed},
    shortcut/.style={->, >=stealth, thin, dotted},
    edgelabel/.style={font=\scriptsize, fill=white, inner sep=0.5pt,
        text height=1.4ex, text depth=0.3ex},
]
  \node[methodnode, centerwidth] (Implementation)  at (5, 7.0) {\ImplementationC};
  \node[methodnode, centerwidth] (MLIPMethod)      at (5, 5.7) {\MLIPMethodC};
  \node[methodnode, centerwidth] (MLIPRun)         at (5, 4.4) {\MLIPRunC};
  \node[methodnode, centerwidth] (TrainedModel)    at (5, 3.1) {\TrainedModelC};
  \node[trainnode,  centerwidth] (MaterialSystem)  at (5, 1.8) {\MaterialSystemC};
  \node[trainnode,  centerwidth] (TrainingDataset) at (5, 0.5) {\TrainingDatasetC};
  \node[methodnode, anchor=west] (Library)              at (-1.1, 7.0) {\LibraryC};
  \node[methodnode, anchor=west] (SimulationType)       at (-1.1, 5.7) {\SimulationTypeC};
  \node[trainnode,  anchor=west] (AtomicConfiguration)  at (-1.1, 1.8) {\AtomicConfigurationC};
  \node[methodnode, minimum width=3.2cm] (HyperparameterSetting) at (11.25, 3.1) {\HyperparameterSettingC};
  \node[benchnode,  anchor=east, minimum width=2.5cm] (BenchmarkStudy)  at (16.5, 3.1) {\BenchmarkStudyC};
  \node[benchnode,  anchor=east, minimum width=2.5cm] (BenchmarkResult) at (16.5, 1.8) {\BenchmarkResultC};
  \node[benchnode,  anchor=east, minimum width=2.5cm] (AccuracyMetric)  at (16.5, 0.5) {\AccuracyMetricC};
  \node[methodnode, minimum width=3.2cm] (FunctionalForm) at (11.25, 7.0) {\FunctionalFormC};
  \node[methodnode, minimum width=3.2cm] (LossFunction)   at (11.25, 5.7) {\LossFunctionC};
  \node[methodnode, minimum width=3.2cm] (Hyperparameter) at (11.25, 4.4) {\HyperparameterC};
  % Leaf classes — single row at y=-0.8
  \node[trainnode]  (DatasetProvenance) at (3,    -0.8) {\DatasetProvenanceC};
  \node[trainnode]  (CoveredProperty)   at (7,    -0.8) {\CoveredPropertyC};
  \node[benchnode]  (MetricType)        at (12.5, -0.8) {\MetricTypeC};
  \node[benchnode, anchor=east]  (MetricProperty)    at (16.5, -0.8) {\MetricPropertyC};

  % DFTSettings: same row as MetricType, centred horizontally in the gap
  % between CoveredProperty's east edge and MetricType's west edge.
  % DFTCalculation is one row above DFTSettings, sharing its x-centre.
  \node[trainnode] (DFTSettings)
    at ($(CoveredProperty.east)!0.5!(MetricType.west)$) {\DFTSettingsC};
  \node[trainnode] (DFTCalculation)
    at ($(CoveredProperty.east)!0.5!(MetricType.west) + (0, 1.3)$) {\DFTCalculationC};

  \draw[edge] (MLIPMethod) -- node[edgelabel] {\hasFunctionalFormR} (FunctionalForm.west);
  \draw[edge] (MLIPMethod) -- node[edgelabel] {\hasLossFunctionR} (LossFunction.west);
  \draw[edge] (MLIPMethod) -- node[edgelabel] {\hasHyperparameterR} (Hyperparameter.west);
  \draw[edge] (MLIPMethod) -- node[edgelabel] {\hasImplementationR} (Implementation);
  \draw[edge] (MLIPMethod) -- node[edgelabel] {\supportsSimulationR} (SimulationType);
  \draw[edge] (Implementation) -- node[edgelabel] {\implementedInR} (Library);
  \draw[edge] (HyperparameterSetting) -- node[edgelabel] {\forHyperparameterR} (Hyperparameter);
  \draw[edge] (MLIPRun) -- node[edgelabel] {\appliesMethodR} (MLIPMethod);
  \draw[edge] (MLIPRun) -- node[edgelabel] {\producesR} (TrainedModel.north);
  \draw[edge] (TrainingDataset) -- node[edgelabel] {\hasDFTCalculationR} (DFTCalculation);
  \draw[edge] (TrainingDataset) -- node[edgelabel] {\coversMaterialR} (MaterialSystem);
  \draw[edge] ($(TrainingDataset.north west)!2/3!(TrainingDataset.south west)$)
    to[out=180, in=270] node[edgelabel, pos=0.6] {\hasConfigurationR} (AtomicConfiguration.south);
  \draw[edge] (DFTCalculation) -- node[edgelabel] {\hasDFTSettingsR} (DFTSettings);
  \draw[edge] (BenchmarkStudy) -- node[edgelabel] {\hasResultR} (BenchmarkResult);
  \draw[edge] (BenchmarkResult) -- node[edgelabel] {\hasAccuracyMetricR} (AccuracyMetric);

  % Leaf-class edges
  \draw[edge] (TrainingDataset) -- node[edgelabel] {\datasetProvenanceR} (DatasetProvenance);
  \draw[edge] (TrainingDataset) -- node[edgelabel] {\coversPropertyR} (CoveredProperty);
  \draw[edge] (AccuracyMetric) -- node[edgelabel] {\metricTypeR} (MetricType);
  \draw[edge] (AccuracyMetric) -- node[edgelabel] {\metricPropertyR} (MetricProperty);

  % runsOn polyline: vertical middle segment runs through the exact
  % midpoint x of the corridor between AtomicConfiguration's east edge
  % and MaterialSystem's west edge.
  \coordinate (runsOnEnd) at ($(TrainingDataset.north west)!1/3!(TrainingDataset.south west)$);
  \coordinate (corridor)  at ($(AtomicConfiguration.east)!0.5!(MaterialSystem.west)$);
  \path let \p1=(corridor), \p2=(MLIPRun.west), \p3=(runsOnEnd) in
    coordinate (runsOnTopCorner) at (\x1, \y2)
    coordinate (runsOnBotCorner) at (\x1, \y3);
  \draw[crossedge, rounded corners=4pt]
    (MLIPRun.west) -- (runsOnTopCorner) --
    node[edgelabel] {\runsOnR}
    (runsOnBotCorner) -- (runsOnEnd);

  % evaluatesModel: leaves BenchmarkResult at 1/4 down its west side and
  % arrives at TrainedModel at 3/4 down its east side; both tangents
  % horizontal-westward (out=180, in=0).
  \draw[crossedge]
    ($(BenchmarkResult.north west)!1/4!(BenchmarkResult.south west)$)
    to[out=180, in=0]
    node[edgelabel, pos=0.5] {\evaluatesModelR}
    ($(TrainedModel.north east)!3/4!(TrainedModel.south east)$);
  \draw[crossedge] (BenchmarkResult) -- node[edgelabel] {\targetMaterialR} (MaterialSystem);

  % HyperparameterSetting incoming: solid from MLIPRun (underlying), dotted
  % shortcut from TrainedModel.  hasHyperparameterSetting leaves MLIPRun
  % heading east (out=0) and arrives at 1/4 down HyperparameterSetting's
  % west side, also heading east (in=180).
  \draw[edge] (MLIPRun.east) to[out=0, in=180]
    node[edgelabel, pos=0.5] {\hasHyperparameterSettingR}
    ($(HyperparameterSetting.north west)!1/4!(HyperparameterSetting.south west)$);
  \draw[shortcut] (TrainedModel.east) to
    node[edgelabel, pos=0.55] {\trainedUsingR}
    (HyperparameterSetting.west);

  % trainedWith shortcut: dotted arc on left of central column, swung
  % wider westward (looseness 2.5) so it does not overlap the runsOn
  % polyline.  Ends at 3/4 down MLIPMethod's west side.
  \draw[shortcut, bend left=70, looseness=2.5]
    (TrainedModel.west) to node[edgelabel, pos=0.5] {\trainedWithR}
    ($(MLIPMethod.north west)!3/4!(MLIPMethod.south west)$);
\end{tikzpicture}
}
\caption{Overview of the \MLIPsOntology{}: the three modules (Algorithm,
  Training Data, Benchmark) and their core classes and relationships.
  Solid nodes are \mlipsprefix{} classes; dashed nodes are external
  superclasses from ML-Schema, PROV-O, and Schema.org.}
\label{fig:ontology-overview}
\end{figure}

\begin{table}[t]
\centering
\caption{Namespace prefixes used by the \MLIPsOntology{}.}
\label{tab:prefixes}
\begin{tabular*}{\textwidth}{l@{\extracolsep{\fill}}l}
\toprule
Prefix & Namespace \\
\midrule
\texttt{mlips:}  & \texttt{\mlipsns} \\
\texttt{mdo:}    & \texttt{https://w3id.org/mdo/core/} \\
\texttt{cmso:}   & \texttt{https://purls.helmholtz-metadaten.de/cmso/} \\
\texttt{asmo:}   & \texttt{https://purls.helmholtz-metadaten.de/asmo/} \\
\texttt{mls:}    & \texttt{http://www.w3.org/ns/mls\#} \\
\texttt{prov:}   & \texttt{http://www.w3.org/ns/prov\#} \\
\texttt{qudt:}   & \texttt{http://qudt.org/schema/qudt/} \\
\texttt{schema:} & \texttt{https://schema.org/} \\
\texttt{wd:}     & \texttt{http://www.wikidata.org/entity/} \\
\bottomrule
\end{tabular*}
\end{table}

\subsection{Algorithm Module}
\label{sec:algorithm-module}

The Algorithm module models MLIP algorithms and their hyperparameter
spaces, addressing CQ1, CQ2, CQ8, and CQ9.

\paragraph{Core classes.}
The central concept is $\MLIPAlgorithmC$, representing a specific
MLIP algorithm (e.g., MACE, ACE, HDNNP)---the architecture itself, not a
trained model. A key defining feature of an MLIP algorithm is
its \emph{atomic environment descriptor}: the mathematical function that
maps a local neighborhood (neighbor positions and types) to an invariant
or equivariant feature vector. Different algorithms use different
descriptors---moment tensors (MTP), SOAP (GAP), symmetry functions
(HDNNP), equivariant message passing (MACE, NequIP), atomic cluster
expansion (ACE)---captured by the $\hasDescriptorR$ role linking
$\MLIPAlgorithmC$ to instances of $\AtomicEnvironmentDescriptorC$.
Each algorithm declares the $\HyperparameterC$ instances it
accepts (cutoff radius, number of layers, learning rate, etc.), with their
names, data types, ranges, and default values. A $\TrainingRunC$ represents
a concrete training execution: it $\executesR$ an $\MLIPAlgorithmC$, takes
a $\TrainingDatasetC$ as input ($\runsOnR$), and $\producesR$ a
$\TrainedModelC$---the learned parameters (weights) that can be used for
prediction. This design follows ML-Schema's activity-based pattern, where
the run mediates between algorithm, data, and model. An $\ImplementationC$
captures a software realization of an algorithm in a specific $\LibraryC$
and version. Finally, $\SimulationTypeC$ represents the kinds of atomistic
simulation an algorithm supports (molecular dynamics, Monte Carlo, geometry
optimization, phonon calculations, thermodynamic integration).
\paragraph{Core properties.}
The role $\hasHyperparameterR$ links an algorithm to the hyperparameters it
accepts, while $\hasImplementationR$ links it to its software
implementations and $\supportsSimulationR$ to supported simulation types.
A $\TrainingRunC$ is connected to its algorithm via $\executesR$, to its
input data via $\runsOnR$, and to the resulting model via $\producesR$.
For convenience, a $\TrainedModelC$ also carries the shortcut roles
$\trainedWithR$ and $\trainedOnR$ (derivable as property chains through the
run), enabling direct queries without navigating the run entity.
Domain-specific datatype properties---$\cutoffRadiusR$, $\numLayersR$,
$\learningRateR$, $\numRadialBasisR$, and $\numAngularBasisR$---record
concrete hyperparameter values, with units annotated via
QUDT~\cite{hodgson2014qudt}. The role $\implementedInR$ connects an
implementation to its library.
Efficiency metadata is split along two dimensions. Prior
knowledge about an algorithm---its asymptotic complexity and whether it
supports GPU or parallel execution---is recorded on $\MLIPAlgorithmC$
via $\dlRole{trainingComplexity}$, $\dlRole{inferenceComplexity}$,
$\dlRole{supportsGPU}$, and $\dlRole{supportsParallelization}$. Measured
training cost (wall-clock time, hardware, peak memory, GPU hours) is
recorded on $\TrainingRunC$, while measured inference cost is recorded
on the corresponding $\BenchmarkResultC$ (see Benchmark Module).

\paragraph{Alignment.}
$\MLIPAlgorithmC$ is a subclass of $\mlsAlgorithmC$, inheriting
ML-Schema's general algorithm description framework. $\TrainingRunC$ is a
subclass of both $\mlsRunC$ and $\provActivityC$, aligning with
ML-Schema's training lifecycle and PROV-O's activity-based provenance model.
$\TrainedModelC$ is a subclass of $\mlsModelC$. $\HyperparameterC$
extends $\mlsHyperParameterC$ with physics-specific properties and
units, and $\ImplementationC$ extends $\mlsImplementationC$.

\paragraph{Axioms.}
The Algorithm module defines 14 axioms: existential restrictions require
every algorithm to declare at least one hyperparameter and simulation type;
cardinality constraints ensure each $\TrainingRunC$ executes exactly one
algorithm, uses one dataset, and produces one model; domain axioms restrict
$\hasHyperparameterR$, $\hasImplementationR$, and $\supportsSimulationR$
to $\MLIPAlgorithmC$. Notably, the shortcut roles on $\TrainedModelC$ are
defined as property chains through the training run:
%
\begin{align*}
  \trainedWithR &\equiv \producesR^{-} \circ \executesR \\
  \trainedOnR &\equiv \producesR^{-} \circ \runsOnR
\end{align*}
%
These allow direct queries (``which algorithm produced this model?'')
without navigating through the $\TrainingRunC$ entity.



\subsection{Training Data Module}
\label{sec:training-data-module}

The Training Data module describes the datasets used to train MLIP models,
addressing CQ3, CQ4, and CQ5.

\paragraph{Core classes.}
A $\TrainingDatasetC$ is a collection of atomic configurations used for MLIP
training. Each dataset covers one or more $\MaterialSystemC$ instances
(elements, alloys, compounds), and its reference data is produced by one or
more $\DFTCalculationC$ instances, each governed by a $\DFTSettingsC$ that
records the exchange-correlation functional ($\xcFunctionalR$, e.g., PBE,
LDA, SCAN), k-point mesh ($\kPointMeshR$), energy cutoff
($\energyCutoffR$), and pseudopotential type ($\pseudopotentialTypeR$).
Individual $\AtomicConfigurationC$ instances represent the atomic positions,
cell parameters, and computed properties within a dataset.
\paragraph{Core properties.}
$\hasDFTSettingsR$ links a calculation to its settings;
$\numConfigurationsR$ records the dataset size;
$\coversPropertyR$ indicates which physical properties
(energy, forces, stresses, virials) are included;
and $\datasetProvenanceR$ classifies the origin
as $\dlConcept{Published}$, $\dlConcept{InHouse}$, or
$\dlConcept{Augmented}$. Augmented datasets are linked to
their sources via PROV-O's \mbox{$\provWasDerivedFromR$}.

\paragraph{Alignment.}
$\TrainingDatasetC$ is modeled as a subclass of both $\mlsDatasetC$
(for ML workflow compatibility) and $\provEntityC$ (for provenance
tracking). $\DFTCalculationC$ aligns with $\mdoCalculationC$.
$\MaterialSystemC$ connects to $\cmsoCrystallineMaterialC$ where
applicable, and $\AtomicConfigurationC$ aligns with
$\cmsoAtomicStructureC$. Material systems link to Wikidata for
chemical element identification.

\paragraph{Axioms.}
The Training Data module defines 5 existential restrictions: every
$\TrainingDatasetC$ must cover at least one material system, declare at
least one covered property, specify its provenance, and include at least
one atomic configuration; every $\DFTCalculationC$ must have DFT settings.
The provenance classification ($\dlConcept{Published}$,
$\dlConcept{InHouse}$, $\dlConcept{Augmented}$) is exhaustive and disjoint.



\subsection{Benchmark Module}
\label{sec:benchmark-module}

The Benchmark module links published studies to trained models, materials,
and accuracy metrics, addressing CQ6 and CQ7.

\paragraph{Core classes.}
A $\BenchmarkStudyC$ represents a published study reporting MLIP evaluation
results. Each study contains one or more $\BenchmarkResultC$ instances, each
capturing a single evaluation: a $\TrainedModelC$ applied to a material
system. Since a trained model already encodes the algorithm, training data,
and hyperparameter settings (via the Algorithm module), the benchmark result
need only reference the model and the test conditions. Results report
$\AccuracyMetricC$ instances (e.g., RMSE of energy, MAE of forces), each
characterized by a $\MetricTypeC$ (RMSE, MAE, R$^2$), a $\MetricPropertyC$
(energy, force, stress), and a numeric value with unit.

\paragraph{Core properties.}
The role $\evaluatesModelR$ links a benchmark result to the trained model
being evaluated, while $\targetMaterialR$ specifies the material system used
for testing. $\hasAccuracyMetricR$ connects to the reported accuracy values.
Each $\AccuracyMetricC$ is described by $\metricTypeR$, $\metricPropertyR$,
and $\metricValueR$. The role $\reportedInR$ links a result to its source
publication.

\paragraph{Alignment.}
$\BenchmarkStudyC$ is modeled as a subclass of $\provActivityC$,
capturing the provenance chain from data to evaluation. Publications link to
$\schemaScholarlyArticleC$ for bibliographic metadata and to Wikidata
for author identification.

\paragraph{Axioms.}
The Benchmark module defines 8 axioms ensuring data completeness: every
$\BenchmarkStudyC$ must contain at least one result; every
$\BenchmarkResultC$ must evaluate exactly one $\TrainedModelC$, target at
least one material system, and report at least one accuracy metric; every
$\AccuracyMetricC$ must specify exactly one metric type and one measured
property. Through the property chains on $\TrainedModelC$, every benchmark
result is transitively linked to the algorithm and training data used,
enabling cross-module queries (CQ6, CQ7) without redundant encoding.



%% ====================================================================
